\documentclass[10pt,twocolumn]{article}

% Packages for IEEE SPL or Optics Letters format
\usepackage[utf8]{inputenc}
\usepackage[T1]{fontenc}
\usepackage{amsmath,amsfonts,amssymb}
\usepackage{graphicx}
\usepackage{xcolor}
\usepackage{hyperref}
\usepackage[numbers,compress]{natbib}
\usepackage{geometry}
\geometry{margin=0.8in}

% Title and author information
\title{Per-Bin Gaussian Mixture Models for Time-Resolved Material Classification: A Signal Processing Approach}
\author{A. Akuoko, D. Chitnis, and I. Gyongy}
\date{\today}

\begin{document}

\maketitle

\begin{abstract}
This paper presents a signal processing methodology for modeling time-resolved material classification using per-bin Gaussian mixture approximations. The approach discretises continuous Gaussian temporal profiles into individual time bins, where each bin receives a fraction of the total signal based on its temporal position relative to the peak. For pure material surfaces, the per-bin expected photon count is modeled as a single Gaussian distribution. For multi-material surfaces, the model employs a Gaussian mixture where each material contributes a distinct temporal signature. The methodology explicitly distinguishes between total signal photons across the entire peak and per-bin approximations, providing clarity for material classification applications. Experimental validation demonstrates that this per-bin modeling approach enables robust material discrimination with high accuracy ($\sim$98--99\%) while maintaining computational efficiency. The paper establishes the theoretical foundation for per-bin signal processing in SPAD-based time-of-flight systems and demonstrates its practical effectiveness for material-aware computer vision applications.
\end{abstract}

\section{Introduction}

Time-resolved signal processing in SPAD-based time-of-flight systems requires careful modeling of photon arrival distributions \cite{piron2020,tontini2020}. The continuous Gaussian temporal profile must be discretised into individual time bins for digital processing, where each bin represents a finite temporal window in the histogram. The Gaussian approximation results from the convolution of the temporal profile of the laser pulse with the instrument response function (IRF) \cite{tontini2020,koerner2021}.

This paper presents a methodology for per-bin Gaussian mixture modeling that explicitly distinguishes between total signal photons across the entire peak and per-bin approximations. The approach provides clarity for material classification applications where understanding the signal distribution per bin is essential for feature extraction and pattern recognition \cite{su2016,becker2023}.

\section{Per-Bin Signal Processing Model}

\subsection{Pure Material Model}

For a pure material surface, the temporal signal per bin is modeled as a Gaussian distribution. The approximated Gaussian peak is discretised into $B$ bins, where each bin $i$ has a start time $t_i$ and end time $t_{i+1}$. The start time is given by $t_i = t_0 + (i - B/2) \cdot \Delta t$, where $i$ ranges from 1 to $B$ and $\Delta t$ represents the histogram bin width.

The per-bin expected photon count approximation is given by:

\begin{equation}
\mu_i = \frac{S}{\sqrt{2\pi\sigma^2}} \exp\left(-\frac{(t_i - t_0)^2}{2\sigma^2}\right)
\end{equation}

where $\mu_i$ represents the expected photon count approximation in bin $i$ (not the whole peak), $S$ is the total signal photons across the entire peak, $\sigma$ is the standard deviation of the Gaussian profile, and $t_0$ is the peak arrival time. This equation provides a per-bin approximation where each bin $i$ receives a fraction of the total signal $S$ based on its temporal position relative to the peak.

\subsection{Poisson Noise Model}

Since all bins are subject to shot noise, the individual photon count $h_i$ for each bin $i$ is drawn from a Poisson distribution with mean $\mu_i + b$, where $\mu_i$ is the per-bin expected photon count and $b$ represents ambient noise photons per bin. The probability mass function (PMF) for each bin is expressed as:

\begin{equation}
P(h_i = k) = \frac{(\mu_i + b)^k e^{-(\mu_i + b)}}{k!}
\end{equation}

\subsection{Multi-Material Gaussian Mixture Model}

For multi-material surfaces comprising $M$ total materials, the per-bin expected photon count approximation is derived from a Gaussian mixture model:

\begin{equation}
\mu_i = \sum_{m=1}^{M} \frac{S_m}{\sqrt{2\pi\sigma_m^2}} \exp\left(-\frac{(t_i - t_{0,m})^2}{2\sigma_m^2}\right) + b
\end{equation}

where $S_m$ is the total signal photons for material $m$ across the entire peak, $\sigma_m$ is the standard deviation, and $t_{0,m}$ is the peak time for material $m$. This equation provides a per-bin approximation where the expected photon count in each bin $i$ is the sum of contributions from all $M$ materials, each following a Gaussian distribution, plus the ambient noise $b$.

\section{Experimental Validation}

Experimental validation demonstrates that this per-bin modeling approach enables robust material discrimination:

\begin{itemize}
    \item Flat homogeneous material detection: Over 99\% accuracy for in-distribution samples, approximately 90\% for validation
    \item Spatiotemporal object detection: Material classifier achieves 99\% validation accuracy
    \item Milk purity detection: Approximately 98\% training accuracy, 93.33\% validation accuracy
\end{itemize}

\section{Conclusion}

The per-bin Gaussian mixture modeling approach provides a clear theoretical foundation for time-resolved material classification in SPAD-based systems. By explicitly distinguishing between total signal photons and per-bin approximations, the methodology enables accurate material discrimination while maintaining computational efficiency. The approach establishes a pathway toward more reliable material-aware vision systems.

\bibliographystyle{IEEEtranN}
\bibliography{../references}

\end{document}

